\paragraph{\trPinStructure。}\label{sub:pin-structures}
$Pin_n$ 是 $O_n$ 经由 $\ZZ/2$ 的一个\textbf{\trCentralExtension}（central
extension）。这样的扩张其实有两个，二者\trTopology{}相同，但作为李群互不相同。就我们的
目的而言，其中一种可能性碰巧比另一种稍微方便一些，于是我们规定：$Pin_n$ 是
（在\trTopologicalGroup{}扩张的意义下）由 $w_2\in H^2(BO_n;\ZZ/2)$ 分类的那个群。
让我们简略回忆一下具体的定义。把 $\RR^n$ \trEmbedding{}到\textbf{\trClifford{}代数}
（Clifford algebra）$Cl(\RR^n)$，在其中向量满足 $v^2=|v|^2e$，这里
$e\in Cl(\RR^n)$ 是单位元。于是 $Pin_n\subset Cl(\RR^n)^\times$ 是由
$S^{n-1}\subset\RR^n$ 生成的乘法子群。\trHomomorphism{} $Pin_n\to O_n$ 是\textbf{\trTwisted
{}伴随作用}（twisted adjoint action）限制在 $\RR^n\subset Cl(\RR^n)$ 上的结果。
这就是说，诸元素 $v\in S^{n-1}\subset Pin_n$ 按 $w\mapsto -vwv$ 作用，这是沿
垂直于 $v$ 的\trHypersurface{}的\trReflection。当然，也可以不用 $\RR^n$，而取
任意一个\trInnerProduct{}空间 $F$ 并定义一个具有同样性质的群 $Pin(F)$。考察诸%
\trReflection{}的\trComposition，便得到：
\begin{lemma}\label{thm:pin-involution-preimage}
设 $A\in O(F)$ 是一个\trInvolution，则 $A$ 在 $Pin(F)$ 中的\trPreimage{}\trCanonical{}地
对应于子空间 $F^{\mathrm{anti}}=\ker(\mathrm{Id}+A)$ 的诸\trOrientation。
任何一个这样的\trPreimage{}的平方是
\begin{equation}\label{eq:pin-preimage-square}
(-1)^{\frac{1}{2}\,\mathrm{rank}(F^{\mathrm{anti}})
(\mathrm{rank}(F^{\mathrm{anti}})-1)}\,e;
\end{equation}
特别地，诸\trReflection{}的\trPreimage{}都是 $2$ 阶元。
\end{lemma}
\noqed
\begin{remark}\label{thm:pin-references}
关于 $Pin$ 的一般文献见 [93]——其中我们的 $Pin_n$ 被记作 $Pin_{0,n}$——以及
[86]，其中它被记作 $Pin_n^+$。另一个版本记作 $Pin_{n,0}$ 或 $Pin_n^-$，它有如下
性质：它关联于类 $w_2+w_1^2$；它含于那个（更常见的）\trClifford{}代数中，其中
向量满足 $v^2=-|v|^2e$；最后，在其中有诸\trReflection{}的\trPreimage{}是 $4$ 阶元。
\end{remark}
如通常那样，一个（实\trInnerProduct）\trVectorBundle{} $F^n\to B$ 上的%
\trPinStructure{}是一个 $Pin_n$-\trPrincipalBundle{} $P^\#\to B$，连同一个%
\trIsomorphism{} $P^\#\times_{Pin_n}\RR^n\cong F$。\trHomomorphism{} $Pin_n\to O_n$ 诱导一个%
\trDoubleCover{} $P^\#\to P$，其中 $P$ 是 $F$ 的\trOrthonormal\trFrame{}丛。由
我们对 $Pin_n$ 的原有刻画，这样一种结构存在的\trObstruction{}是 $w_2(F)$。若它
为零，则诸 $Pin_n$-结构在相差\trIsomorphism{}的意义下的可能选取构成
$H^1(B;\ZZ/2)$ 上的一个仿射空间。具体地说，给定一个 \trPinStructure{}
$P^\#$ 和一个\trRealLineBundle{} $\beta\to B$，可以用\trFibreProduct{}
$P^\#\times_{\ZZ/2}S(\beta)$ 作出一个新的 \trPinStructure：其中 $\ZZ/2$ 在关联于
$\beta$ 的\trDoubleCover{} $S(\beta)$ 上以 $\pm1$ 作用，而在 $P^\#$ 上则通过中心
子群 $\ker(Pin_n\to O_n)=\{\pm e\}$ 作用。为简便计，我们将把新结构记作
$P^\#\otimes\beta$。
对任意\trVectorBundle{} $F$，都可以关联一个群丛 $Pin(F)$。我们想澄清一个容易
引起混淆之点。给定 $F$ 上的一个 \trPinStructure{} $P^\#$，考虑群丛
$\mathrm{Aut}(P^\#)\to B$，其在 $x$ 上的\trFiber{}由\trPrincipalHomogeneousSpace
{} $Pin_n$-空间 $P_x^\#$ 的全部\trAutomorphism{}组成。这可以看作一个相伴丛
\begin{equation*}
\mathrm{Aut}(P^\#)=P^\#\times_{Pin_n}Pin_n=P\times_{O_n}Pin_n
\hookrightarrow P\times_{O_n}Cl(\RR^n).
\end{equation*}
这里 $Pin_n$ 通过共轭作用于自身；这下降为一个 $O_n$-作用，后者继而\trEmbedding{}到
$Cl(\RR^n)$ 上的非\trTwisted{}伴随作用。例如，沿\trHyperplane{} $v^\perp$ 的\trReflection
{}按 $w\mapsto vwv$ 作用，而限制到 $Cl(\RR^n)$ 中的 $\RR^n\subset Cl(\RR^n)$ 上，
它是该\trReflection{}的 $-1$ 倍。由此可见，$\mathrm{Aut}(P^\#)$ 不%
\trIsomorphic{}于 $Pin(F)$，而是\trIsomorphic{}于
$Pin(F\otimes\lambda^{\mathrm{top}}(F))$；并且这纳入下图
\begin{equation}\label{eq:aut-pin-diagram}
\begin{tikzcd}
\mathrm{Aut}(P^\#)\ar[r,"\cong"]\ar[d]
  & Pin(F\otimes\lambda^{\mathrm{top}}(F))\ar[d] \\
\mathrm{Aut}(P)=O(F)\ar[r,"A\mapsto A\otimes\mathrm{Id}","\cong"']
  & O(F\otimes\lambda^{\mathrm{top}}(F)).
\end{tikzcd}
\end{equation}
更一般地，若取一个 \trPinStructure{} $P^\#$ 并把它与两个\trLineBundle{}
$\beta_0$、$\beta_1$ 作张量，则所得的两个\trPrincipalBundle{}之间的
$Pin_n$-\trEquivariant{}映射丛是
\begin{equation}\label{eq:pin-equivariant-iso}
\mathrm{Iso}(P^\#\otimes\beta_0,\ P^\#\otimes\beta_1)\cong
Pin(F\otimes\lambda^{\mathrm{top}}(F))\times_{\ZZ/2}
S(\beta_0^\vee\otimes\beta_1).
\end{equation}
\begin{remark}\label{thm:twisted-pin-structures}
还有更一般的\textbf{\trTwisted{} \trPinStructure}（twisted Pin structure）
理论：相对于某个参考类 $w\in H^2(B;\ZZ/2)$ 而\trTwisted{}的 \trPinStructure。其存在的%
\trObstruction{}是 $w_2(F)+w$。为给出一个 \trGerbe{}式的定义，选取分别表示 $w$
与丛 $F$ 的 Cech 上链 $\check{w}\in\check{C}^2(B;\ZZ/2)$ 与
$\check{F}\in\check{C}^1(B;O_n)$。于是，一个\trTwisted{} \trPinStructure{}就是
$\check{C}^1(B;Pin_n)$ 的一个元素，它\trProjection{}到 $\check{F}$ 上，而它未能成为%
\trCocycle{}这件事则由 $\check{w}$ 编码（由于 $\ZZ/2\subset Pin_n$ 是一个中心
子群，这是有意义的）。
\end{remark}
另一种图像不那么一般，但就前面讨论过的材料而言或许更容易把握，具体如下。设有
一条\trSmooth\trHypersurface{} $\Delta\subset B$，其上带一个\trRealLineBundle{}
$\beta$，而 $w$ 由 $w_1(\beta)\in H^1(\Delta;\ZZ/2)$ 经\trPushforward{}得到。为
使叙述简明，我们还假设 $\Delta$ 是\trCoorientation{}的。沿 $\Delta$ 剖开 $B$
得到一个\trManifold{} $\bar{B}$，它有两条\trBoundaryComponent{}
$\partial_+\bar{B}$、$\partial_-\bar{B}\cong\Delta$。于是，$F$ 上的一个%
\trTwisted{} \trPinStructure{}是\trPullback{} $\widetilde{F}\to\bar{B}$ 上的一个%
\trPinStructure{} $\widetilde{P}^\#$，连同\trIsomorphism{}
$\widetilde{P}^\#|_{\partial_+\bar{B}}\cong
\widetilde{P}^\#|_{\partial_-\bar{B}}\otimes\beta$，其\trCovering{}显然的等同
$\widetilde{F}|_{\partial_-\bar{B}}=F|_\Delta=\widetilde{F}|_{\partial_+\bar{B}}$。
$Pin$ 之所以与 $\bar\partial$-算子的\trIndex{}理论相关，是它与 $w_2$ 之间关系的
一个后果。为看到一个基本的例子，我们简短地回到
\cref{thm:det-line-bundle-loop-space}与\cref{thm:abstract-brane-loop-determinant}
的情形。考虑一条\trLoop{} $\rho\in\scr{L}_k\,Gr(V)$、相伴的\trVectorBundle{}
$F_\rho\to S^1$，以及\trDisk{}上的算子 $D_{D,\rho}$。设 $\widetilde{F}_\rho$ 是
$F_\rho$ 在把两端点等同起来的映射 $[0;1]\to S^1$ 下的\trPullback。我们想给它
配备
\begin{equation}\label{eq:twisted-pin-structure}
\text{\emph{一个 \trPinStructure{} $\widetilde{P}_1^\#$，以及一个\trIsomorphism}}
\qquad
\widetilde{P}_1^\#\cong\widetilde{P}_0^\#\otimes
\lambda^{\mathrm{top}}(\rho(1))^{\otimes k}.
\end{equation}
当 $k$ 为偶数时，这就是 $F_\rho$ 上的一个 \trPinStructure。当 $k$ 为奇数时，它
是\cref{thm:twisted-pin-structures}意义下\trTwisted{} \trPinStructure{}定义的
一个特例，其中：$F=F_\rho$；\trHypersurface{}是点 $\Delta=\{0\}\subset S^1$；
而 $\beta\to\{0\}$ 无非是实直线 $\lambda^{\mathrm{top}}(\rho(1))$。这显然可以
推广到\trLoop{}的\trFamily，此时\eqref{eq:twisted-pin-structure} 的存在性%
\trObstruction{}来自类
\begin{equation}\label{eq:twisted-pin-obstruction}
T(w_2)+k\,U(\mu)\in H^1(\scr{L}_k\,Gr(V);\ZZ/2).
\end{equation}
\begin{lemma}\label{thm:twisted-pin-determinant-iso}
在 $F_\rho$ 上选取\trTwisted{} \trPinStructure{}
\eqref{eq:twisted-pin-structure}，便确定一个\trIsomorphism{}
$\det(D_{D,\rho})\cong\lambda^{\mathrm{top}}(\rho(0))$。
\end{lemma}
\begin{proof}
对每个 $\rho$，\trTwisted{} \trPinStructure{}都有两个可能的\trIsomorphism{}类，
它们构成 $\scr{L}_k\,Gr(V)$ 的一个\trDoubleCover，由
\eqref{eq:twisted-pin-obstruction} 分类。由
\cref{thm:det-line-bundle-loop-space}，它\trIsomorphic{}于由以
$\det(D_{S,\rho})^\vee\otimes\lambda^{\mathrm{top}}(\rho(0))$ 为\trFiber{}的%
\trLineBundle{}给出的那个\trCovering。要在这两个丛之间确定一个\trIsomorphism，
因而只需考虑它们在 $\scr{L}_k\,Gr(V)$ 的某个非空\trConnected{}子集上的限制。
这与\cref{thm:abstract-brane-loop-determinant}所用的策略相当类似；事实上 $k=0$
的情形只不过是那个证明的重新表述。考察常值\trLoop{} $\rho$：此时丛 $F_\rho$ 是
平凡的，其上遂有一个特异的平凡 \trPinStructure。规则是：若该%
\trPinStructure{}平凡，则选取由 $\ker(D_{D,\rho})\cong\rho(0)$ 诱导的%
\trIsomorphism{} $\det(D_{D,\rho})\cong\lambda^{\mathrm{top}}(\rho(0))$；若它
非平凡，则改变其符号。
现在设 $k=-1$。取一个\trOrthogonal{}复\trSplitting{} $V=V_+\oplus V_-$，其中
$V_-$ 的复维数为一，并对某个\trSplitting\trLagrange{}子空间
$\Lambda=\Lambda_+\oplus\Lambda_-$ 考虑形如
$\rho(s)=\Lambda_+\oplus\exp(-\pi s I_V)\Lambda_-$ 的\trLoop{}。$F_\rho$ 上有一个
显然的\trConnection，其\trMonodromy{}为
$A=\mathrm{Id}_{\Lambda_+}\oplus(-\mathrm{Id}_{\Lambda_-})$。\trTwisted{}
\trPinStructure{} \eqref{eq:twisted-pin-structure} 的一个选取由把 $A$
\trLift{}为诸\trPinStructure{}之间的一个\trIsomorphism{}给出：
$\widetilde{P}_0^\#\otimes\lambda^{\mathrm{top}}(\rho(0))\to\widetilde{P}_1^\#$。
施用\eqref{eq:pin-equivariant-iso} 可以发现，这对应于
$Pin(\Lambda\otimes\lambda^{\mathrm{top}}(\Lambda))\times_{\ZZ/2}
S(\lambda^{\mathrm{top}}(\Lambda))$ 中 $A\otimes\mathrm{Id}$ 的一个\trPreimage。由
\cref{thm:pin-involution-preimage}，两个 $\lambda^{\mathrm{top}}(\Lambda)$ 相互
抵消，于是这化归为 $\Lambda_-$ 的一个\trOrientation。另一方面，把算子%
\trSplitting{}成 $V_\pm$ 两块并施用
\cref{thm:det-line-bundle-loop-space}证明的第 (iii) 部分，可以看出 $D_{D,\rho}$
是\trSurjective{}的，且其\trKernel{}\trIsomorphic{}地映到 $\Lambda_+$ 上。于是，
$\Lambda_-$ 的一个\trOrientation{}\trCanonical{}地确定一个\trIsomorphism{}
$\det(D_{D,\rho})\cong\lambda^{\mathrm{top}}(\Lambda)$，正如所论。这关于
$V=V_+\oplus V_-$ 的\trSplitting{}与 $\Lambda=\Lambda_+\oplus\Lambda_-$ 的选取
都是连续的。
对一般的 $k$，考虑形如 $\rho(s)=\Lambda_+\oplus\exp(k\pi s I_V)\Lambda_-$ 的%
\trLoop{}。论证中有关 \trPinStructure{}的部分只依赖于 $k$ 的奇偶性，故与从前相同。
在\trDeterminantLine{}丛一侧，则可利用\trGluing{}公式
\eqref{eq:boundary-gluing} 得到\trCanonical{}的\trIsomorphism{}
$\det(D_{D,\rho})\cong\lambda^{\mathrm{top}}(\Lambda_+)\otimes
(\Lambda_-)^{\otimes(1+k)}$。
\end{proof}
